\chapter{多項式}
    \section{規約}
        \begin{itemize}
            \item 以降の説明に於いて多項式中に特別の断りなく現れる$x$は不定元を意味する。
            \item 体$F$上の、$x$を不定元とする多項式全体の集合を$F[x]$と表記する。
            \item 体$F$上の、$x$を不定元とするLaurent多項式全体の集合を$F[x,x^{-1}]$と表記する。
        \end{itemize}
    \section{係数が次数に関して対称な奇数次の多項式\texorpdfstring{$f(x)$}{f(x)}を\texorpdfstring{$(x+1)$}{x+1}で除したものも係数が次数に関して対称になる}
        \begin{proof}
            \quad\par
            $n\in\naturalNumbers$を奇数とし、$f(x) = \sum_{i=0}^n a_ix^i,\;a_i = a_{n-i}$とする。
            $f(x)$は$(x+1)$で割り切れる。
            $g(x) = f(x)/(x+1) = \sum_{i=0}^{n-1} b_ix^i$としたときに$b_i = b_{n-1-i}$となることを示す。
            \begin{align*}
                f(x) &= (x+1)g(x) = \sum_{i=0}^{n-1} (b_ix^{i+1} + b_ix^i) = b_0 + \sum_{i=1}^{n-1}b_ix^i + \sum_{i=1}^n b_{i-1}x^i \\
                &= b_0 + \sum_{i=1}^{n-1} (b_{i-1} + b_i)x^i + b_{n-1}x^n
            \end{align*}
            これよりまず$b_0 = a_0 = a_n = b_{n-1}$となる。
            次に$x,x^2,\dotsc,x^{n-1}$の項の係数に関しては$b_{i-1} + b_i = a_i = a_{n-i} = b_{n-1-i} + b_{n-i}$すなわち$b_i - b_{n-1-i} = b_{(n-1)-(i-1)} - b_{i-1}$となる。
            既に見た$b_0 = b_{n-1}$から始めて帰納的に示せる。
        \end{proof}
    \section{\texorpdfstring{$\sum_{i=0}^n a_ix^i$}{}が既約多項式ならば\texorpdfstring{$\sum_{i=0}^n a_{n-i}x^i$}{}も既約多項式である}
        \label{既約多項式の係数の順番反転}
        \begin{shadebox}
            体$F$上の$n$次多項式$f_1(x) \coloneqq \sum_{i=0}^n a_ix^i$が$F$上で既約であれば、係数の順番を逆にした多項式$f_2(x) \coloneqq \sum_{i=0}^n a_{n-i}x^i$もまた既約である。
        \end{shadebox}
        \begin{proof}
            \quad\par
            背理法で示す。
            $f_2(x)$が既約でないと仮定すると、適当な多項式$g_{2,1}(x),g_{2,2}(x) \in F[x],\;d_{2,1} + d_{2,2} \coloneqq \deg(g_{2,1}(x)) + \deg(g_{2,1}(x)) = n$が存在して次式が成り立つ。
            \[ f_2(x) = g_{2,1}(x)g_{2,2}(x) \]
            上式をLaurent多項式に拡張して両辺の$x$を$x^{-1}$で置き換えると
            \[ f_2(x^{-1}) = g_{2,1}(x^{-1})g_{2,2}(x^{-1}) \]
            両辺に$x^n$を掛けると
            \[ f_1(x) = x^n f_2(x^{-1}) = x^{d_{2,1}}g_{2,1}(x^{-1})x^{d_{2,2}}g_{2,2}(x^{-1}) =: g_{1,1}(x)g_{1,2}(x) \]
            ここに$g_{1,1}(x), g_{1,2}(x)$は$g_{2,1}(x), g_{2,2}(x)$の係数の順番を逆にしたものである。
            これは$f_1(x)$が既約であることに矛盾する。
        \end{proof}
    \section{既約多項式\texorpdfstring{$\iff$}{⇔}最小多項式}
        \begin{shadebox}
            $P$を体$F$上の多項式とする。
            $P$が$F$上で既約であることと、$P$がその任意の根の最小多項式であることは同値である。
        \end{shadebox}
        \begin{proof}
            \quad\newline\noindent
            ($\Rightarrow$)
            \par
            背理法で示す。
            $\alpha$を$P$の任意の根とする。
            仮に$P$より低次の$F$上の多項式$P_2$が存在して$P_2(\alpha)=0$であるとする。
            $P$を$P_2$で割った商を$Q_2$, 余りを$P_3$とすると、$0 = P(\alpha) = Q_2(\alpha)P_2(\alpha) + P_3(\alpha)$より$P_3(\alpha) = 0$である。
            同じ要領で$P_i$を$P_{i+1}$で割った商を$Q_{i+1}$, 余りを$P_{i+2}$とすると$0 = P(\alpha) = P_2(\alpha) = P_3(\alpha) = \dotsb$となる。
            $P_i$の次数は$i$の増加と共に真に減少するため、次のいずれかが成り立つ。
            \begin{enumerate}
                \item ある$i$で$P_{i+1} \mid P_i$
                \item ある$i$で$\deg(P_i)=0$となる
            \end{enumerate}
            Euclidの互除法により、1.の場合は$\gcd(P,P_2) = \gcd(P_2,P_3) = \dotsb = \gcd(P_i,P_{i+1}) = P_{i+1}$となり、$P$が規約であるという前提に矛盾する。
            2.の場合も先述の通り$P_i(\alpha)=0$であり、$\deg(P_i)=0$に注意して$P_i=0$であるので$P_{i-1} \mid P_{i-2}$であり、1.と同様に矛盾が生じる。
            \newline\newline
            ($\Leftarrow$)
            \par
            背理法で示す。
            $\alpha$を$P$の任意の根とする。
            仮に$P$が既約でないとすると$F$上の$P$より低次な2つの多項式$P_2,P_3$を用いて$P=P_2P_3$と表せる。
            $\alpha$は$P$の根であるから$0 = P(\alpha) = P_2(\alpha)P_3(\alpha)$より$P_2(\alpha) = 0$または$P_3(\alpha) = 0$であるが、これは$P$が$\alpha$の最小多項式であることに矛盾する。
        \end{proof}